% FILENAME: paper/section-Cosmology.tex

\section{Cosmological Bounds and Compatibility Limits}
% TODO for Author: Provide a brief introductory paragraph stating that this section verifies the framework's compatibility with established physical limits (non-relativistic mechanics, cosmology, and asymptotic limits) without needing to rely on ad-hoc parameters.

\subsection{The Schrödinger Non-Relativistic Limit}
% TODO for Author: Explain that as a compatibility check, splitting the large and small components of the covariant Dirac field natively recovers the standard 1/2m non-relativistic Schrödinger Hamiltonian.

\begin{equation}
    2m \Psi_{\text{small}} = P_+ (\partial_0 \Psi + m \Psi) + \gamma^i P_- (\partial_i \Psi)
\end{equation}
\begin{minted}{lean}
-- CGD.Quantum.Schroedinger
theorem exactSchroedingerReduction (dPsi : Fin 4 → SpacetimePoint → Matrix (Fin 4) (Fin 4) Complex) 
  (Psi : SpacetimePoint → Matrix (Fin 4) (Fin 4) Complex) (m : Complex) (x : SpacetimePoint) :
  localDiracOp (fun a => dPsi a x) = m • Psi x →
  let D0_mod := modulatedTemporalDeriv (dPsi 0 x) (Psi x) m;
  let Psi_small := P_plus * Psi x;
  (2 • m • Psi_small = P_plus * D0_mod + gammaVec 1 * (P_minus * dPsi 1 x) + gammaVec 2 * (P_minus * dPsi 2 x) + gammaVec 3 * (P_minus * dPsi 3 x)) ∧
  (P_minus * D0_mod = gammaVec 1 * (P_plus * dPsi 1 x) + gammaVec 2 * (P_plus * dPsi 2 x) + gammaVec 3 * (P_plus * dPsi 3 x))
\end{minted}

\subsection{Non-Singular Cosmological Origins}
% TODO for Author: Explain how enforcing a topological boundary condition means the Big Bang manifests as a pure Euclidean SO(4) instanton bounce rather than a mathematical singularity. Additionally, explain how parity inversion naturally flips the topological charge, linking the arrow of time to matter/antimatter asymmetry.

\begin{equation}
    g_{\mu\nu} = c \delta_{\mu\nu}
\end{equation}
\begin{minted}{lean}
-- CGD.Cosmology.BigBang
theorem kinematicBigBang (u : Universe) (phaseRegion : Set SpacetimePoint)
  [vol : CGD.Axioms.MacroscopicVolume u phaseRegion]
  (h_symm : ∀ x ∈ phaseRegion, isFully4DSymmetric (fun mu nu => curvatureSl2c u.sd_sector mu nu x)) :
  ∀ x ∈ phaseRegion, ∃ c : Complex, c ≠ 0 ∧ urbantkeMetric (fun mu nu => curvatureSl2c u.sd_sector mu nu x) = c • 1
\end{minted}

\begin{equation}
    F_{0i} \to -F_{0i} \implies Q \to -Q
\end{equation}
\begin{minted}{lean}
-- CGD.Cosmology.ParityInversion
theorem kinematicParityInversion (u : Universe) :
  ∀ (x : SpacetimePoint) (P_F : Fin 4 → Fin 4 → SL2C),
    isParityInvertedTensor (fun m n => curvatureSl2c u.sd_sector m n x) P_F x →
    pontryaginDensity P_F = - pontryaginDensity (fun mu nu => curvatureSl2c u.sd_sector mu nu x)
\end{minted}

\subsection{Asymptotic String Breaking}
% TODO for Author: Briefly mention that while the 1D flux tube represents a confining Strong Force potential, topological crossover limits dictate that the string will eventually mathematically "snap" at macroscopic distances. 

\begin{equation}
    L > \frac{2M}{\sigma} \implies \text{String Breaking}
\end{equation}
\begin{minted}{lean}
-- CGD.Quantum.Entanglement.Decay
theorem algebraicStringBreakingLimit
  (energyFunc : (Fin 4 → SpacetimePoint → SL2C) → ℝ)
  (intactState snappedState : ℝ → Fin 4 → SpacetimePoint → SL2C)
  {sigma M : ℝ} [eb : FluxTubeStringBreaking (Fin 4 → SpacetimePoint → SL2C) energyFunc intactState snappedState sigma M]
  (u : Universe) (L : ℝ)
  (h_L_pos : L > 0)
  (h_intact : u.sd_sector = intactState L)
  (h_L_crossover : L > (2 * M) / sigma) :
  ¬ isGlobalMinimum energyFunc u.sd_sector
\end{minted}
